\documentclass[11pt]{article}
\usepackage[utf8]{inputenc}
\usepackage[T1]{fontenc}
\usepackage{lmodern}
\usepackage[margin=1in]{geometry}
\usepackage{amsmath,amssymb}
\usepackage{textcomp}
\usepackage[dvipsnames]{xcolor}
\usepackage{enumitem}
\usepackage{parskip}
\usepackage{titlesec}
\usepackage[colorlinks=true,linkcolor=blue,citecolor=blue,urlcolor=blue]{hyperref}

\titleformat{\section}{\large\bfseries}{\thesection}{1em}{}
\titleformat{\subsection}{\normalsize\bfseries}{\thesubsection}{1em}{}
\titleformat{\subsubsection}{\normalsize\bfseries\itshape}{}{0em}{}

\title{\textbf{Agreement, Not Aggregate:\\ Measuring Where Automated Constraint Verification Matches Human Review}}
\author{}
\date{}

\newcommand{\ph}{\textit{[placeholder --- fill in after experiments are run]}}

\begin{document}
\maketitle
\vspace{-2em}

\begin{quote}
\small\textit{Assembled from the six confirmed stage outputs of an \texttt{auto-research-discussion} run (seed topic T5.3), post research-taste audit (\textcolor{ForestGreen}{\textbf{PASS}}). Status: \textbf{specification / paper draft --- no experiments have been run.} All numbers below are pre-registered design targets, not results.}
\end{quote}

\begin{abstract}
Coding agents now produce patches faster than teams can manually check them against explicit design constraints, and automated constraint verification is increasingly treated as a safe way to close that gap --- but whether an automated verifier's precision and recall actually match human review, and for which kinds of constraints, has never been measured against a human-agreement reference rather than an assumed-perfect one. We build a three-way-aligned evaluation: category-labeled design constraint cards, coding-agent-produced patches drawn from real issue-to-pull-request tasks, and adjudicated multi-reviewer judgments over the same pairs. A hybrid verifier --- deterministic rule and architecture-conformance checks for structurally-defined constraints, a calibrated, abstention-capable judge for intent-level ones --- is scored per constraint category against the adjudicated reference, and against the human reviewers' own agreement with each other, rather than against an assumed-infallible ground truth. Categories are classified automatable, human-required, or reference-unstable from that comparison, and a budget-matched single-judge control isolates whether any observed split reflects the constraints themselves or the verifier's own construction. We report where automated constraint checking already matches human review, where it does not, and where the human reference itself is too noisy to judge against at all.
\end{abstract}

\section{Introduction}

Coding agents are now routinely trusted to produce patches against explicit design constraints --- architectural boundaries a module must not cross, APIs a patch must not call, naming conventions a codebase enforces, security invariants that must hold, design principles like separation of concerns that a card can name but not fully spell out. As the volume of agent-produced patches grows past what any team can manually check against every constraint, a natural response has taken hold: build an automated verifier, point it at the constraint cards and the patches, and let it flag violations before a human ever has to look. This is cheap to build and easy to justify, which is exactly why it is worth checking rather than assuming.

Yet a basic question about this practice remains unmeasured: does an automated verifier's precision and recall at catching constraint violations actually match human review, and does that hold uniformly across constraint types, or only for some? Treating a forbidden-API call and a violation of ``preserve separation of concerns'' as the same kind of detection problem, and reporting one blended accuracy number across both, hides exactly the information a team would need to decide where automated checking is safe to rely on and where a human still has to look. Worse, human review itself is being treated implicitly as an infallible reference in this practice --- but if reviewers disagree with each other about whether a card was violated, an apparent verifier failure could really be unresolved disagreement about what the card means, a distinction that determines whether the fix belongs in the verifier or in the card.

The nearest existing work to this question already puts all three ingredients --- explicit design constraints, coding-agent patches, and an automated compliance checker --- on the table: ``Does Pass Rate Tell the Whole Story?'' (arXiv:2604.05955, April 2026) mines and validates 1{,}787 design constraints from real pull requests across six SWE-bench-Verified/Pro-aligned repositories and scores agent patches for compliance with an LLM-based verifier. We concede that overlap plainly, because the delta is exactly where this paper lives: that benchmark \emph{trusts} its verifier as the scoring instrument in order to measure agents, whereas this paper makes the verifier itself the measured object --- asking, per constraint category, whether any such automated checker's judgments actually agree with human review, and whether human review even agrees with itself enough to be the reference. A benchmark built on an LLM verifier inherits precisely the unmeasured quantity this study is designed to measure. A second adjacent line, automated code review built around defect-focused merge-request analysis, balances precision and coverage when flagging bugs against a defect oracle, and reports that fewer than one in ten automatically generated review comments on prior benchmarks are truly high quality --- direct evidence that automated review signal quality is not solved even before constraint-specific structure enters the picture; but it is measured against a bug-defect oracle, not a taxonomy of explicit design constraints scored per category. Remaining adjacent lines each supply a different piece without closing the gap directly. Selective-prediction frameworks for LLM judges bound aggregate disagreement with human labels via a calibrated confidence threshold, but do not decompose where that disagreement concentrates by constraint type. Checklist- and multi-constraint-following work shows that per-constraint verifiability already varies by constraint type --- the empirical premise this paper's category split rests on --- but studies whether a \emph{generator} satisfies constraints, not whether a \emph{verifier} correctly detects violations against human-labeled ground truth. Execution-verified and execution-free software-engineering benchmarks establish realistic, repository-scale coding-agent patches and reward signals for whole trajectories, but their oracles either cover only the executably-testable slice of constraints or score trajectory quality rather than per-constraint human agreement. And verifier-capability studies on math and reasoning tasks show that verification quality does not scale simply with model size and varies strongly by regime --- a warning this paper takes seriously --- but those studies use automatic, not human, ground truth, leaving the human-review noise ceiling itself unmeasured in any domain resembling this one.

Turning this into a measurable question surfaces its own difficulties. Any verifier's performance report is only as informative as its granularity: an aggregate accuracy number can hide a verifier that handles syntactic constraints well and semantic ones poorly behind a blended score that looks acceptable overall, so the category-level split itself has to be measured with enough statistical power to be trusted, not just asserted. The human reference against which the verifier is judged is not guaranteed to be reliable either --- reviewers can disagree with each other about whether a patch violates a card, and unless that disagreement is measured and reported as an explicit ceiling, a verifier's apparent failure cannot be distinguished from the reference's own unresolved noise. A patch can also satisfy a constraint's literal wording while violating what the card actually intends, or vice versa, which no verifier relying only on surface pattern-matching against a card's text can detect. And precision and recall trade off differently by constraint type --- tuning for recall on one category can inflate false positives past the point that drives real-world tool abandonment, while the same threshold may be too permissive on another category entirely --- so no single fixed operating point can be adoptable across every constraint type at once.

Our approach turns each of these difficulties into a specific, checkable design choice. Every verifier score is reported per constraint category rather than pooled into one number, so hidden variation across categories becomes visible by construction rather than by luck. Human review's own reliability is measured directly --- a multi-reviewer panel's agreement with itself is computed and reported as an explicit ceiling, and the verifier's agreement with the adjudicated reference is judged relative to that ceiling, not against an assumed-perfect one; where the ceiling itself is too low to be usable, the category is reported as reference-unstable rather than silently scored anyway. Intent-laden constraints are routed to a judge that can abstain rather than guess when its confidence is low, while structurally-defined constraints are checked deterministically --- and a budget-matched control with no such routing isolates whether any observed category split reflects the constraints' own nature or an artifact of how the verifier is built. Rather than a single precision/recall threshold, a per-category operating curve is reported, so the trade-off that drives real-world adoption is visible category by category. And the underlying corpus of patches is deliberately pre-screened for non-obviousness, so that a verifier's apparent success is not simply an artifact of a corpus where every violation was easy to spot.

We evaluate this design on constructed, category-labeled design cards paired with coding-agent patches seeded from real issue-to-pull-request tasks, scored against a naive keyword matcher and a budget-matched single-judge control, with all comparisons made per constraint category against a measured, not assumed, human-agreement ceiling.

In summary, we contribute:
\begin{itemize}[leftmargin=1.4em]
\item A category-stratified evaluation protocol that reports an automated constraint verifier's precision, recall, and agreement with human review separately per constraint category, rather than as one blended number that can hide exactly where automation fails.
\item A measured, rather than assumed, human-agreement ceiling --- computed directly from a multi-reviewer panel --- against which verifier performance is judged, together with an explicit reference-unstable outcome for categories where that ceiling itself is too low to license any comparison.
\item An operationalizability-routed verifier design, paired with a budget-matched single-judge control, that isolates whether an observed automatable/human-required split reflects the constraints' own nature or the verifier's own construction.
\item A pre-screened, non-obvious three-way-aligned corpus --- design cards, coding-agent patches, and adjudicated human labels --- and a per-category operating-curve reporting protocol that together make the precision/recall trade-off driving real-world adoption visible category by category, rather than hidden behind a single threshold.
\end{itemize}

\section{Problem Definition}

\noindent\textbf{Formal setup.} Let $C$ be a set of design constraint cards, each assigned by a category map $\tau: C \to T$ to a category in $T = \{\text{structural}, \text{API}, \text{naming}, \text{security}, \text{intent-level}\}$, where the reliability of this assignment is itself checked (inter-annotator agreement) rather than assumed. Let $P$ be a set of coding-agent-produced patches, with the evaluation unit a pair $(c,p) \in C \times P$ where card $c$ is applicable to patch $p$. An automated verifier $V$ maps each pair to a judgment $V(c,p) \in \{\text{violation}, \text{no-violation}\}$ with an associated confidence $q(c,p) \in [0,1]$. A panel of $r \geq 3$ independent human reviewers supplies labels $h_1(c,p),\dots,h_r(c,p) \in \{\text{violation}, \text{no-violation}\}$, adjudicated into a reference label $y(c,p)$. For each category $t \in T$, let $\kappa_{H,t}$ denote the human--human agreement (Fleiss' $\kappa$) among the panel, $\kappa_{V,t}$ denote the verifier--human agreement (Cohen's $\kappa$) between $V$ and $y$, and $n_t$ denote the number of observed violations in category $t$. The object of study is a function from $(\kappa_{H,t}, \kappa_{V,t}, n_t)$ to one of four category-level labels --- automatable, human-required, reference-unstable, or inconclusive --- computed separately for every category, never as one pooled statistic.

\medskip
\noindent\textbf{Estimand.} The quantity of interest is the per-category gap $\Delta_t = \kappa_{H,t} - \kappa_{V,t}$ between measured human--human and verifier--human agreement, with unit of analysis one adjudicated (card, patch) pair; the reference population is category-labeled design-card $\times$ coding-agent-patch pairs constructible from real issue-to-pull-request tasks; and the decision this estimand feeds is a binary per-category delegation choice --- route that constraint category to automated verification, or keep it under mandatory human review --- never a paper-wide verdict on ``the verifier.''

\noindent\textbf{Why this is hard.} Fixing the category map $\tau$ does not fix a uniformly reliable verification target: categories differ intrinsically in how directly their truth condition can be read off a patch's observable structure versus how much it depends on design intent that a card names but does not encode, so a verifier tuned to succeed on one category is not guaranteed to succeed on another, and an aggregate statistic over all of $C$ would hide exactly this variation. The reference $y$ itself is not guaranteed to be reliable --- $\kappa_{H,t}$ can be low for categories where reviewers themselves disagree about what a card requires, in which case comparing $V$ against $y$ at all risks attributing the reference's own noise to the verifier. Even where $y$ is reliable, $V$'s judgment can fail in a specific way that a purely structural check cannot: a patch can satisfy a card's literal text while violating its design intent, or vice versa, which requires reasoning about intent rather than pattern-matching against wording. And because the acceptable balance of false positives against false negatives is not the same across categories --- tuning for recall on one category can push false positives past the threshold that drives reviewers to abandon a tool, while the same operating point may be too conservative for another category --- no single confidence threshold can be assumed adoptable across all of $T$ at once.

\section{Method Overview}

At a conceptual level, the study is a single measurement design with five cooperating parts. Design cards and coding-agent patches enter the pipeline already organized by category; a scoring layer keeps every verifier statistic separated by that category rather than pooling them; a human panel supplies not just a reference label but its own measured reliability; a routed verifier construction handles structurally-defined and intent-laden categories by different mechanisms, with a control that checks whether that routing itself is shaping the results; and a corpus-construction step ensures the patches being scored are neither trivially compliant nor trivially violating in the first place.

The \textbf{category-stratified scorer} takes the verifier's per-pair judgments and computes precision, recall, F1, and agreement statistics separately for each of the five categories, refusing to collapse them into one aggregate figure. The \textbf{measured human-agreement ceiling} takes the raw, pre-adjudication labels from the human panel and computes how much reviewers agree with each other, per category, establishing an empirical ceiling rather than assuming the adjudicated reference is infallible. The \textbf{operationalizability-routed verifier} takes each (card, patch) pair and applies a deterministic rule or architecture-conformance check to structurally-defined categories, and a calibrated, abstention-capable judge to intent-laden ones, while a budget-matched control without this routing tests whether the resulting category-level pattern is intrinsic to the constraints or an artifact of the routing itself. The \textbf{operating-curve reporting} layer takes the scored judgments and reports a confidence-thresholded precision/recall curve per category rather than a single blended number. The \textbf{pre-screened corpus construction} step governs which (card, patch) pairs enter the study at all, discarding pairs the human panel finds trivially obvious and requiring a minimum number of retained violations before a category is included in the analysis.

\section{Methodology}

\subsection{Category-Stratified Scorer}
\textbf{Challenge.} An aggregate accuracy figure can hide a verifier that performs well on syntactically checkable constraints and poorly on semantic ones behind a blended score that looks acceptable overall.

\textbf{Insight.} Reporting-granularity, not verifier engineering, is what surfaces hidden category-level variation --- the variation already exists in the data whether or not it is measured, so ensuring it is visible is a design choice about how results are computed, not about the verifier itself.

\textbf{Method.} Precision, recall, F1, and Cohen's $\kappa$ against the adjudicated reference are computed independently for each of the five constraint categories. Every automatable/human-required/reference-unstable/inconclusive verdict is issued at the category level; no pooled, cross-category statistic is reported as a headline result. Because every such verdict is scored per category, the category taxonomy is itself treated as a measurement instrument and validated before any scoring: each category is grounded in a named external source (architecture-conformance literature for structural constraints, static-analysis rule taxonomies for API-usage and naming, security-invariant checklists for security, design-principle literature for intent-level); category boundaries are refined on a held-out pilot labeling round, with any category pair falling below the agreement threshold merged or re-drawn via a named operation rather than silently re-labeled; assignment reliability is validated with at least two independent annotators at $\kappa \geq 0.7$; the annotators' confusion structure is inspected, so that a systematic collapse of one category into another is reported as an ill-defined taxonomy rather than scored; and the category set is frozen before corpus scale-up.

\subsection{Measured Human-Agreement Ceiling}
\textbf{Challenge.} Judging a verifier against human review as if human review were an infallible reference conflates a real verifier failure with the reviewers' own unresolved disagreement about what a card requires.

\textbf{Insight.} Since human review is the yardstick the verifier is measured against, the yardstick's own reliability has to be measured directly, not assumed --- a verifier can only meaningfully be said to fall short of a ceiling that is itself known to be attainable.

\textbf{Method.} A panel of at least three independent reviewers labels every (card, patch) pair before adjudication, and Fleiss' $\kappa$ among the raw reviewer labels is computed per category as an explicit ceiling. The verifier's Cohen's $\kappa$ against the adjudicated reference is compared to that ceiling via a bootstrapped confidence interval on their gap, rather than compared to an assumed-perfect reference; categories whose ceiling itself falls below a usability floor are reported as reference-unstable rather than scored as if the reference were trustworthy.

\subsection{Operationalizability-Routed Verifier with Calibrated Abstention}
\textbf{Challenge.} A patch can satisfy a constraint card's literal wording while violating its underlying design intent, or the reverse --- a failure mode that surface pattern-matching against a card's text cannot detect, and that is concentrated in intent-laden categories that have no reliable per-item checkable predicate.

\textbf{Insight.} Constraints that are not reducible to a deterministic check should not be forced through one, and a verifier that must answer every pair with equal confidence is forced to guess exactly where its judgment is least informative; abstention converts a forced wrong guess into an honest deferral.

\textbf{Method.} Structurally-defined categories are checked with a deterministic rule or architecture-conformance mechanism. Intent-laden categories are routed to an LLM-as-judge mechanism that reasons over a card's stated text together with its design rationale, calibrated to abstain --- deferring to human review --- below a stated confidence threshold. A budget-matched control, using a single-shot judge with no card decomposition, no category routing, and no calibrated abstention, is run at the same total token spend per pair as the full routed verifier; comparing the two isolates whether an observed category-level pattern reflects the constraints' own operationalizability or the routing mechanism itself.

\subsection{Per-Category Operating-Curve Reporting}
\textbf{Challenge.} Precision and recall trade off at a different rate for every constraint category, so a single fixed confidence threshold applied uniformly is, by construction, wrong for at least some categories.

\textbf{Insight.} The reportable unit for adoption decisions is a curve, not a point --- a blended, category-agnostic threshold hides exactly the trade-off information a team needs to decide whether a verifier is usable for a given constraint type at all.

\textbf{Method.} A confidence-thresholded precision/recall operating curve is computed per category, alongside a comparison against a naive keyword/regex matcher that uses no structural or semantic reasoning. Categories where the naive matcher already falls within the human-agreement band are reported as not requiring further verifier engineering, independent of whether the full verifier also outperforms it there.

\subsection{Pre-Screened Corpus Construction}
\textbf{Challenge.} A corpus of (card, patch) pairs that is trivially easy --- where every violation is obvious at a glance --- would make any verifier configuration, including a naive keyword matcher, look accurate, hiding the real automatable/human-required split behind an artifact of corpus difficulty.

\textbf{Insight.} Class balance between violating and compliant pairs does not guarantee genuine difficulty; a pair can be balanced yet trivial, so difficulty has to be pre-screened for directly rather than inferred from counts alone.

\textbf{Method.} Pairs are seeded from real issue-to-pull-request coding-agent tasks rather than synthetic templates. The human panel pre-screens candidate pairs and discards those judged trivially compliant or trivially violating, measured by whether the reviewer's judgment time for that pair falls below a non-obviousness threshold. A category enters the main analysis only once it retains a minimum share of non-obvious pairs and a minimum number of adjudicated violations; categories that do not clear this bar are reported as inconclusive rather than folded into an aggregate.

\section{Experiments}

\subsection{Experimental Setup}

\textbf{Datasets.} A constructed design-card corpus, category-labeled across the five constraint categories and decomposed into item-level checkable predicates --- the necessary substrate for any category-stratified result, since without category-labeled cards there is no unit to stratify by. Coding-agent patches seeded from real issue-to-pull-request tasks (SWE-bench / SWE-bench Verified-style) --- used specifically because ecologically real patches are necessary for the non-obviousness pre-screen to mean anything; a fully synthetic corpus would let violations be constructed as obvious by design, defeating the pre-screen's purpose. Constructed violating-patch variants, added per category where real agent patches naturally yield too few violations to retain a category in the analysis --- used to avoid conflating corpus scarcity with a null finding on constraint operationalizability. Where category coverage permits, cards may additionally be seeded from the mined-and-validated constraint inventory of the design-aware issue-resolution benchmark (arXiv:2604.05955), which supplies pull-request-provenanced constraints already aligned to SWE-bench-Verified/Pro instances; any card so seeded is re-labeled under this paper's five-category taxonomy and passes the same inter-annotator reliability gate as constructed cards.

\medskip
\noindent\textbf{Annotation budget (pre-registered design target).} The study is annotation-dominated, so the budget is stated up front rather than discovered late: at roughly 120 retained (card, patch) pairs per category $\times$ 5 categories $\approx$ 600 pairs, a $\geq 3$-reviewer panel at an estimated 6 minutes per pair-label implies $\approx$ 180 reviewer-hours for the main labeling pass, plus $\approx$ 40 hours for the non-obviousness pre-screen (lighter, discard-oriented judgments over a larger candidate pool), $\approx$ 30 hours for adjudication of disagreements, and $\approx$ 20 hours for the double-coded error-taxonomy pass --- $\approx$ 270 reviewer-hours total, the scale of a funded annotation effort rather than a weekend pilot. The walking-skeleton phase (two categories, API-usage and intent-level, at $\approx$ 40 pairs each) costs $\approx$ 15\% of this and is the kill-test: if the pilot's Fleiss' $\kappa$ or retained-violation counts make the full design's power targets unreachable, the full-scale pass is not run.

\textbf{Baselines} (sanity-check and rival-explanation comparators, not competitors to beat --- this is a measurement study, and no research question here is framed as ``our verifier beats baseline X'' as its headline claim). A naive keyword/regex matcher, using no AST parsing and no LLM reasoning, establishes the floor a verifier needs to clear to be worth building at all --- and, more informatively, reveals which categories already need no further verifier research because the naive matcher already reaches the human-agreement band there. A single-shot LLM judge with no card decomposition, no category routing, and no calibrated abstention, budget-matched in total token spend per pair to the full routed verifier, isolates whether any observed category-level pattern is intrinsic to the constraints or an artifact of the routing/decomposition mechanism. A calibrated selective-judge protocol in the style of prior LLM-judge selective-prediction work is reported descriptively alongside these two, as the nearest published calibrated-judge design, with no pass/fail criterion attached to it.

\textbf{Metrics.} Per-category precision, recall, and F1 against the adjudicated reference; per-category Cohen's $\kappa$ (verifier--human) alongside Fleiss' $\kappa$ (human--human, the measured ceiling); the bootstrapped confidence interval on their gap; and the resulting per-category classification (automatable / human-required / reference-unstable / inconclusive).

\textbf{Research questions:}
\begin{enumerate}[leftmargin=1.6em]
\item RQ1 --- Does the verifier reach a reliable precision/recall bar on structurally-defined constraint categories?
\item RQ2 --- Is the intent-level category's performance significantly worse than the structural categories', not just numerically lower?
\item RQ3 --- Is the human-review reference itself usable --- do reviewers agree with each other enough to serve as a ceiling?
\item RQ4 --- Which categories reach the human-agreement band (automatable), which fall short (human-required), and which are inconclusive?
\item RQ5 --- Does an observed category-level performance split reflect the constraints' own nature, or the verifier's routing mechanism?
\item RQ6 --- For which categories does even a naive keyword matcher already suffice, and does the full verifier improve on it elsewhere?
\item RQ7 --- Does a single confidence threshold work across all constraint categories, or does the adoptable operating point vary by category?
\item RQ8 --- Is the evaluation corpus non-trivial enough, per category, to support a real finding rather than an artifact of easy cases?
\end{enumerate}

\subsubsection{RQ1: Does the verifier reach a reliable precision/recall bar on structurally-defined constraint categories?}
\textbf{Setup.} For each of the four structurally-defined categories (architectural boundary, API-usage restriction, naming/style, security invariant) separately, compute precision and recall against the adjudicated reference with Wilson 95\% confidence intervals. Sample size: each category retains at least 30 adjudicated violations (per the corpus pre-screen) as its inclusion floor; the interval-width arithmetic below is stated against that floor rather than hand-waved.

\textbf{Expected outcome.} Two-tier, matched to what the sample size can actually support. Tier 1 (all four structural categories, at the 30-violation floor): precision $\geq 0.85$ and recall $\geq 0.80$ per category, with Wilson 95\% CI lower bounds $\geq 0.60$ for both. Tier 2 (only categories retaining $\geq 80$ adjudicated violations): CI lower bounds $\geq 0.75$. The tiering is forced by Wilson arithmetic, not preference: at $n = 30$, an observed recall of 0.80 yields a lower bound of $\approx 0.63$ and even 0.90 yields $\approx 0.74$, so a 0.75 lower-bound target at the 30-violation floor would be refuted by sample size alone; at $n = 80$ an observed recall of 0.85 yields $\approx 0.75$.

\textbf{Falsification condition.} Refuted if any structural category's point estimate falls below its precision/recall bar, or its CI lower bound falls below the tier bound its sample size registers it for --- that would mean even the most operationalizable constraint categories are not yet reliably automatable at prototype maturity, a materially different finding from the intent-level gap this study otherwise expects.

\textbf{Hypothesis (plain language).} Structurally-defined constraints --- the ones whose truth condition is computable from a patch's observable structure --- are where an automated verifier should already be closest to human-level reliability; if it is not, that is itself the headline finding.

\textbf{Metrics.} Per-category precision, recall, and their Wilson 95\% confidence intervals.

\textbf{Result.} \ph

\subsubsection{RQ2: Is the intent-level category's performance significantly worse than the structural categories', not just numerically lower?}
\textbf{Setup.} Fit a logistic regression of per-pair verifier correctness (correct/incorrect against the adjudicated reference) on category, with an agent covariate and card as a clustering unit for standard errors; compare the intent-level category's F1 to the minimum F1 among the four structural categories.

\textbf{Expected outcome.} The intent-level category's F1 falls at least 0.15 below the minimum structural-category F1, and the category effect in the regression is statistically significant via a likelihood-ratio test at $\alpha = 0.05$.

\textbf{Falsification condition.} Refuted if the F1 gap is smaller than 0.15, or if the category effect is not statistically significant --- either would mean the intent/structural divide this study is built around is not actually borne out as a distinct, significant effect, as opposed to ordinary noise across categories.

\textbf{Hypothesis (plain language).} The gap between structural and intent-level constraint categories is not just a numerically lower score but a real, statistically distinguishable effect of category.

\textbf{Metrics.} Per-category F1; likelihood-ratio test statistic and p-value on the category coefficient.

\textbf{Result.} \ph

\subsubsection{RQ3: Is the human-review reference itself usable --- do reviewers agree with each other enough to serve as a ceiling?}
\textbf{Setup.} Compute Fleiss' $\kappa$ among the raw, pre-adjudication labels from the $r \geq 3$-reviewer panel, separately for each of the five categories.

\textbf{Expected outcome.} Fleiss' $\kappa$ is at least 0.4 in every category, establishing the human--human reference as usable rather than merely assumed usable.

\textbf{Falsification condition.} Refuted for any category where Fleiss' $\kappa$ falls below 0.4 --- that category is then reported as reference-unstable, and any verifier comparison against it in RQ1/RQ2/RQ4 is treated as uninterpretable rather than as evidence about the verifier.

\textbf{Hypothesis (plain language).} Before judging the verifier against human review, the study has to establish that human review agrees with itself enough to be a meaningful yardstick in the first place.

\textbf{Metrics.} Per-category Fleiss' $\kappa$.

\textbf{Result.} \ph

\subsubsection{RQ4: Which categories reach the human-agreement band (automatable), which fall short (human-required), and which are inconclusive?}
\textbf{Setup.} For each category with a usable reference (RQ3), compute the gap $\Delta_t$ between the human--human ceiling $\kappa_{H,t}$ and the verifier--human agreement $\kappa_{V,t}$, with a 95\% bootstrap confidence interval (10,000 resamples, clustered by card), applying a Holm correction across the five categories for the human-required tests; categories with fewer than 30 observed violations are forced to inconclusive regardless of the interval (per RQ8's gate).

\textbf{Expected outcome.} At least two categories are classified automatable ($\Delta_t$'s CI upper bound $\leq 0.10$), the intent-level category is classified human-required ($\Delta_t$'s CI lower bound $> 0.10$), and at most one category is classified inconclusive.

\textbf{Falsification condition.} Refuted if fewer than two categories reach the automatable band, or if the intent-level category's CI does not clear the human-required bar, or if more than one category is inconclusive --- any of these would mean the presumed structural/semantic divide does not cleanly produce the expected automatable/human-required partition.

\textbf{Hypothesis (plain language).} The automatable/human-required split is not all-or-nothing --- most, but not all, structurally-defined categories should reach a human-comparable agreement band, while the intent-level category should not.

\textbf{Metrics.} Per-category $\Delta_t$ and its bootstrap confidence interval; resulting three-way classification.

\textbf{Result.} \ph

\subsubsection{RQ5: Does an observed category-level performance split reflect the constraints' own nature, or the verifier's routing mechanism?}
\textbf{Setup.} Run the budget-matched single-shot judge (no card decomposition, no routing, no abstention) on the same corpus, and compare its per-category F1 pattern to the full routed verifier's pattern. Because ``the patterns agree'' is an equivalence claim, it is tested as one: per category, a 95\% bootstrap confidence interval (clustered by card) is computed on the F1 difference between the two configurations, and equivalence is claimed via two one-sided tests (TOST) against a pre-registered margin of $\pm 0.05$ --- never by eyeballing point estimates, whose per-category sampling noise at the 30--80-violation scale can alone exceed 0.05.

\textbf{Expected outcome.} For every category, the bootstrap CI on the F1 difference lies entirely within $\pm 0.05$ (TOST equivalence holds), indicating that any observed automatable/human-required split reflects the constraints' intrinsic operationalizability rather than the routing or decomposition mechanism.

\textbf{Falsification condition.} Refuted if any category's difference CI lies entirely outside the $\pm 0.05$ margin --- that would mean the routing/decomposition design itself is shaping which categories look automatable, and any automatable/human-required claim would have to be reported as instrument-conditional rather than about the constraints themselves. If a category's CI straddles the margin (neither equivalence nor difference established), that category's RQ5 verdict is reported as underpowered-inconclusive rather than resolved in either direction.

\textbf{Hypothesis (plain language).} If a category looks automatable or human-required, that should be a property of the constraint type, not an artifact of how the verifier happens to be built.

\textbf{Metrics.} Per-category F1 for both verifier configurations; the per-category F1 difference with its clustered-bootstrap 95\% CI and TOST equivalence verdict against the $\pm 0.05$ margin.

\textbf{Result.} \ph

\subsubsection{RQ6: For which categories does even a naive keyword matcher already suffice, and does the full verifier improve on it elsewhere?}
\textbf{Setup.} Score the naive keyword/regex matcher on the same corpus; compare its per-category agreement with the adjudicated reference against the human-agreement band from RQ4, and compare its paired decisions against the full verifier's via McNemar's test ($\alpha = 0.05$, Holm-corrected across the five categories).

\textbf{Expected outcome.} On any category where the naive matcher already falls inside the human-agreement band, that category is reported as not requiring further verifier engineering; on remaining categories, the full verifier's paired decisions improve on the naive matcher's at a statistically significant level.

\textbf{Falsification condition.} Refuted (in the direction that would matter) if the full verifier fails to improve on the naive matcher, at a significant level, on any category outside the naive matcher's own already-sufficient band --- that would mean the added routing/LLM-judge machinery is not earning its complexity on the categories that actually need it.

\textbf{Hypothesis (plain language).} The informative question is not whether a sophisticated verifier beats a naive one in aggregate --- a hybrid verifier beating a keyword matcher overall would not be surprising --- but which specific categories the naive matcher already handles well enough that no further verifier research is needed there, and whether the added machinery earns its keep on the rest.

\textbf{Metrics.} Per-category agreement of the naive matcher against the human-agreement band; McNemar's test statistic and Holm-corrected p-value per category.

\textbf{Result.} \ph

\subsubsection{RQ7: Does a single confidence threshold work across all constraint categories, or does the adoptable operating point vary by category?}
\textbf{Setup.} Compute a confidence-thresholded precision/recall operating curve per category; identify, per category, the lowest confidence threshold at which precision reaches 0.85.

\textbf{Expected outcome.} No single confidence threshold clears a 0.85 precision floor simultaneously on all five categories --- the precision-maximizing threshold differs significantly across categories.

\textbf{Falsification condition.} Refuted if a single, category-invariant confidence threshold clears the 0.85 precision floor on all five categories simultaneously --- that would mean a uniform operating point is adoptable everywhere, contrary to the per-category trade-off this design assumes.

\textbf{Hypothesis (plain language).} Because precision and recall trade off differently by constraint type, no single deployment threshold should be adoptable across every category at once --- the operating point itself is a per-category design decision.

\textbf{Metrics.} Per-category precision/recall operating curves; per-category threshold at which precision first reaches 0.85.

\textbf{Result.} \ph

\subsubsection{RQ8: Is the evaluation corpus non-trivial enough, per category, to support a real finding rather than an artifact of easy cases?}
\textbf{Setup.} Among pairs surviving the human pre-screen, compute the share per category whose median reviewer judgment time exceeds the pre-specified non-obviousness threshold, and the retained adjudicated violation count per category.

\textbf{Expected outcome.} At least 20\% of surviving pairs per category are judged non-obvious, and at least 30 adjudicated violations are retained per category, before that category is included in RQ1--RQ7's analyses.

\textbf{Falsification condition.} Refuted for any category failing either bar --- that category is reported as inconclusive across all other RQs regardless of what its raw precision/recall numbers would otherwise show, rather than folded into the main findings on an under-powered or artificially easy basis.

\textbf{Hypothesis (plain language).} A result is only as trustworthy as the corpus it was measured on; a category with too few real violations or only obvious ones cannot support a finding, no matter how clean its raw numbers look.

\textbf{Metrics.} Per-category share of non-obvious pairs; per-category retained violation count.

\textbf{Result.} \ph

\vspace{0.5em}
\noindent\textbf{Mechanism analysis (supplementary).} Beyond establishing \emph{that} and \emph{where} a category-level gap exists, a pre-registered error-taxonomy pass explains \emph{why}: on every verifier error the adjudicated panel disagrees with, a second independent coding pass (two raters, disagreements adjudicated) classifies the error as (i) a literal-vs-intent mismatch, (ii) a structural/pattern-matching failure, or (iii) other. The budget covers all errors on the two walking-skeleton categories (API-usage and intent-level) plus a stratified sample of at least 20 errors per category at full scale, reported as an error-taxonomy table (category $\times$ mechanism-type counts) alongside the headline $\Delta_t$ numbers. The pre-committed reading: if literal-vs-intent mismatches dominate intent-level errors while structural/pattern-matching failures dominate structural-category errors, the observed gap is mechanistically confirmed as caused by the literal/intent distinction; if structural failures dominate intent-level errors instead, the gap is an engineering artifact of the judge component rather than evidence that intent-level constraints are inherently human-required.

\medskip
\noindent\textbf{Further evaluation.} With more time, a natural extension is cross-agent and cross-repository generalization --- testing whether the automatable/human-required partition found here holds for other coding agents, languages, and codebases, rather than the fixed agent/repository set used here. This is explicitly out of scope for the present design: it would require a second corpus-construction and pre-screening pass rather than a re-analysis of the existing one, and the present study is scoped to establish whether the three-way-alignment protocol itself produces a usable, non-trivial reading at all before that generalization question is worth asking.

\section{Related Work}

The closest prior work to this paper is the design-aware issue-resolution benchmark of ``Does Pass Rate Tell the Whole Story?'' (arXiv:2604.05955, April 2026), and it is worth stating the delta plainly before surveying the wider landscape: that work mines and validates explicit design constraints from real pull requests (495 issues, 1{,}787 constraints, six SWE-bench-Verified/Pro-aligned repositories) and checks agent-patch compliance automatically with an LLM-based verifier --- every substrate ingredient of the present study. The difference is the direction of measurement: there, the verifier is the trusted instrument and the agents are the measured objects; here, the verifier is the measured object and adjudicated multi-reviewer human judgment --- with its own agreement measured, per category, as an explicit ceiling --- is the reference. Neither the per-category verifier-vs-human agreement gap, the human--human ceiling, nor a reference-unstable outcome appears in that benchmark's design, which is precisely the gap this paper is built to close; conversely, its mined-and-validated constraint inventory is a natural seeding source for this paper's card corpus.

\textbf{Automated code review and LLM-judge calibration.} Defect-focused automated code review reframes review around merge requests and key-bug detection, combining static code slicing for dependency-aware context with a multi-role pipeline that explicitly balances precision against coverage, and reports that fewer than 10\% of automatically generated review comments on prior benchmarks are truly high quality. Selective-prediction frameworks for LLM judges answer only when an estimated human-agreement confidence clears a calibrated threshold, escalating uncertain cases through a cascade of stronger judges to bound aggregate disagreement risk. Both establish that automated review and automated judgment carry a real, non-trivial precision cost, but neither decomposes that cost by an explicit constraint taxonomy --- the defect-focused work has no notion of constraint category at all, and the selective-judge framework bounds disagreement in aggregate without asking which categories drive the abstentions.

\textbf{Constraint decomposition and instruction-following.} Checklist-based alignment work decomposes an instruction into atomic, separately-checkable yes/no requirements, arguing that global rubrics miss instruction-specific requirements --- a template this paper's card decomposition borrows directly. Complex-instruction-following benchmarks synthesize multi-constraint instructions paired with validation mechanisms and show that constraint satisfaction degrades as constraint count grows, establishing that per-constraint verifiability already varies by type. Both provide the empirical premise behind an automatable/human-required split, but both study whether a \emph{generator} satisfies constraints, not whether a \emph{verifier} correctly detects human-labeled violations in already-produced patches --- the detection-precision question this paper measures directly.

\textbf{Software-engineering agent evaluation.} Execution-verified benchmarks pair real GitHub issues with resolving pull requests, using fail-to-pass tests as an objective correctness oracle over repository-scale edits, and are the source of realistic coding-agent patches this paper's corpus draws on. Execution-free reward models score whole agent trajectories instead, motivated by test suites being incomplete or misaligned with what actually matters. Design-aware issue resolution (arXiv:2604.05955) extends this line past the execution oracle by scoring the same repository-scale patches against mined design constraints via an LLM verifier --- but, as stated above, treats that verifier as trusted rather than as an object of measurement. In the adjacent reward-hacking setting, EvilGenie (arXiv:2511.21654) does verify its automated detectors --- an LLM judge, held-out tests, and edit detection --- against human review, finding the judge effective in unambiguous cases; that validation is aggregate and domain-specific, with no constraint taxonomy and no measurement of how much human reviewers agree with each other. None of these measures how well an automated checker matches human review on a fixed, explicit constraint per patch, per constraint category, against a human--human agreement ceiling --- the execution oracle only covers the executably-testable slice of constraints, the trajectory-level reward model was never built to isolate a single card's compliance, and the design-aware benchmark's verifier is exactly the instrument whose reliability this paper measures.

\textbf{Verifier-capability measurement, architecture conformance, and tool adoption.} Verification-dynamics analysis empirically disentangles problem difficulty, generator capability, and verifier capability on math and reasoning tasks, showing verifier quality does not scale simply with model size and cautioning that a single aggregate accuracy number hides real per-regime variation --- the direct motivation for this paper's category-stratified reporting. Architecture-conformance techniques (Reflexion models) compute convergences and divergences between an intended high-level design model and as-built source structure, the canonical example of a structurally checkable, automatically verifiable design constraint, though the underlying tool predates coding agents by two decades and offers no evidence on semantic or intent-level constraints. Interview-based studies of static-analysis tool adoption find that false positives, not false negatives, are the dominant reason developers abandon automated checkers in practice --- establishing that precision, not just recall, governs real-world usability, a finding this paper's per-category operating-curve reporting is built to make actionable rather than merely cite.

\section{Conclusion}

This work turns a governance assumption already embedded in how coding-agent output is checked --- that an automated verifier can be trusted to catch design-constraint violations --- into a measured claim with an explicit reference. The central methodological commitment is that a verifier's agreement with human review is only interpretable relative to how much human reviewers agree with each other in the first place, and only meaningful when reported per constraint category rather than as one blended figure; a study that skipped either commitment could not distinguish a genuine automatable category from a lucky aggregate, or a genuine human-required category from unresolved reviewer noise being misattributed to the verifier. Whichever categories turn out to be automatable, the finding is decision-relevant in a direct, practical sense: it tells a team exactly where an automated constraint gate can be trusted and where human review remains necessary, rather than leaving that decision to an untested assumption. Two limitations bound this claim from the outset. The design is scoped to a fixed set of coding agents, repositories, and a five-category taxonomy validated for inter-annotator agreement; whether the same automatable/human-required partition holds for other agents, languages, or a finer-grained taxonomy is a natural next step rather than a claim made here. And the reference-unstable outcome this design can produce is itself informative --- a category where human reviewers cannot agree with each other is not a verifier failure to fix, but a card-design problem to address before automation can even be meaningfully evaluated there, pointing toward clearer card authoring as a companion research direction rather than more verifier engineering alone.

\vspace{1em}
\hrule
\vspace{1em}

\section*{Pre-handoff checklist}
\begin{itemize}[leftmargin=1.4em]
\item \textbf{Organization/navigation} --- narrative flows problem $\to$ method $\to$ experiments $\to$ related work without forward references to undefined terms.
\item \textbf{Terminology discipline} --- ``constraint category,'' ``human-agreement ceiling,'' ``reference-unstable,'' ``operating curve'' each defined at first use in Problem Definition, used consistently thereafter; no internal stage IDs (\texttt{sp\_*}, \texttt{m\_*}, \texttt{rq\_*}) leaked into prose --- checked directly against the discussion pipeline log.
\item \textbf{Notation with worked example} --- the formal setup section defines $(C, \tau, T, P, V, q, h_1,\dots,h_r, y, \kappa_{H,t}, \kappa_{V,t}, n_t)$ and the Problem Definition/Methodology sections consistently reuse this notation; no separate worked numeric example is included since no numbers exist yet (would require fabricating results).
\item \textbf{RQ framing matches outcomes} --- all 8 RQs are graded/quantitative (rates, CIs, F1 gaps, thresholds), each answered in its own subsection with an explicit falsifiable expected outcome and a distinct falsification condition, not a binary yes/no.
\item \textbf{No typesetting artifacts} --- only explicit placeholder markers in Results, per the drafting convention this is a pre-registered spec, not a report.
\item \textbf{Artifact plan} --- the design-card corpus, patch corpus, adjudicated human labels, and verifier prototype are the natural release artifacts; not yet built, so not claimed as available.
\end{itemize}

\section*{Hand-off}
This is a draft, not a finished paper --- experiments have not run. Natural next steps: prose polish and venue-specific LaTeX formatting once a target venue is chosen, or revising an earlier stage if piloting surfaces a problem (e.g., if Fleiss' $\kappa$ falls below the 0.4 usability floor in more categories than expected).

\end{document}
